\par\noindent\textbf{(1)} 令结界面x=0，耗尽区−xp<x<0有ρ=−qNA，0<x<xn有ρ=+qND，外区ρ≈0。积分Poisson dE/dx=ρ/ε，并令边缘E=0，界面场连续要求qNA xp=qND xn，因此xp/xn=ND/NA，体现离子空间电荷总量抵消。
\par\noindent\textbf{(2)} 电场在两侧为线性三角形，峰值绝对值|E\_max|=qNAxp/ε=qNDxn/ε。总势降ΔV=∫|E|dx=|E\_max|W/2，W=xp+xn；同时ΔV=V\_bi+V\_R，因此由W和浓度的代数关系求空间电荷区域。
\par\noindent\textbf{(3)} 将xp=ND W/(NA+ND)、xn=NA W/(NA+ND)代入，得W²=(2ε(V\_bi+V\_R)/q)(1/NA+1/ND)。单位面积耗尽结电容C\_A=dQ\_A/dV\_R的幅值为ε/W，随反偏压按(V\_bi+V\_R)\textasciicircum{}\{-1/2\}下降；正向小信号还有扩散电容，不可一概忽略。
\par\noindent\textbf{(4)} 在高反偏下W增长、|E\_max|增长，若电场达到雪崩或隧穿击穿阈值，静态耗尽近似和恒定离子分布不能再描述电流。低偏和高注入时中性区少子扩散与复合主导，结电容包含扩散电容，不再只由ε/W给出。
